\subsection*{A.1. Quantitation of Perfluorohexanesulfonamide in Plasma}

Quantification of perfluorohexanesulfonamide (PFHxSAm) in plasma samples was completed by RTI International (Research Triangle, NC). An ultra-performance liquid chromatography-tandem mass spectrometry (UPLC-MS/MS) method was developed to determine PFHxSAm concentrations in rat plasma. A six-point matrix calibration curve, in the range of 10--1,000~ng/mL, was prepared by adding 10~$\mu$L of an appropriate spiking standard (PFHxSAm in methanol) and 10~$\mu$L of an internal standard solution (perfluoro-1-[$^{13}$C$_{8}$]octanesulfonamide in methanol) to 100~$\mu$L of control matrix (adult male Sprague Dawley rat plasma). Quality control (QC) samples were prepared similarly at a target concentration of 50~ng/mL in plasma. Blanks and study samples were prepared like calibration standards, except 10~$\mu$L of methanol was used in place of spiking solution. Sample extraction was achieved by adding 300~$\mu$L of acetonitrile, vortex mixing, and centrifuging for 10 minutes. A 10~$\mu$L aliquot from the supernatant of each sample was then diluted with 990~$\mu$L of 80/20 acetonitrile/water and mixed.

All samples were analyzed using a Waters Acquity UPLC (Milford, MA) coupled with an Applied Biosystems API-5000 mass detector (Framingham, MA). A Waters Acquity UPLC HSS T3 column (2.1~$\times$~100~mm, 1.8~$\mu$m) and HSS T3 guard (2.1~$\times$~5~mm, 1.8~$\mu$m) were used with mobile phases A (0.1\% formic acid in water) and B (0.1\% formic acid in acetonitrile). A flow rate of 0.3~mL/min was used with 1\% B for 1 minute, a linear gradient of 1\% to 99\% B in 15 minutes, held for 4 minutes, reversed to 1\% B in 0.5 minutes, and held for 1.5 minutes. The electrospray ion source was operated in the negative ion mode with a source temperature of 650$^{\circ}$C and an ion spray voltage of $-4500$~V. Transition ranges monitored were m/z 398.0 to 78.0 (PFHxSAm quantitation ion) and m/z 506.0 to 78.0 (internal standard).

A linear regression with 1/X$^{2}$ weighting was used to relate peak area ratios of analyte to internal standard and analyte concentrations. Calibration curves were linear ($r > 0.99$). The lower limit of quantitation (LOQ) for PFHxSAm in rat plasma was 10.0~ng/mL. For QC samples, the accuracy measured as percent relative error was within $\pm16.0\%$ of the nominal concentration. The concentrations (ng/mL) of PFHxSAm in study samples were calculated using peak area ratios and the regression equation. Samples that had concentrations above the upper limit of the calibration curve were diluted and reanalyzed. All values were above the LOQ and were reported.
